% =============================================================================
% Appendix G: Battery Adapter Interface Specification
% ORIUS/DC3S Battery-Safety Thesis
% =============================================================================
\chapter{Battery Adapter Interface Specification}
\label{app:adapter-interface}

This appendix specifies the \texttt{BatteryDomainAdapter} interface that
bridges the domain-agnostic DC3S pipeline to the battery energy storage
domain.  Every domain that instantiates ORIUS must implement these six
methods; the battery adapter serves as the reference implementation.

\section{Interface Overview}

The adapter exposes six methods that map one-to-one onto the DC3S
pipeline stages: Observe, Qualify, Inflate, Tighten, Repair, Certify.
Table~\ref{tab:adapter-methods} summarises each method.

\begin{table}[htbp]
\centering
\caption{\texttt{BatteryDomainAdapter} interface methods.}
\label{tab:adapter-methods}
\small
\begin{tabular}{@{}p{3.0cm}p{3.0cm}p{3.0cm}p{4.0cm}@{}}
\toprule
\textbf{Method} & \textbf{Input} & \textbf{Output} & \textbf{Battery interpretation} \\
\midrule
\texttt{ingest\_telemetry}
  & Raw BMS packet $\ztt$
  & Parsed state vector $\ot$
  & Extracts SOC, voltage, current, temperature from the degraded telemetry stream.
\\[4pt]
\texttt{compute\_oqe}
  & $\ot$, history buffer
  & Reliability score $\wt \in [0,1]$
  & Aggregates freshness, consistency, and drift evidence into a scalar quality score.
\\[4pt]
\texttt{build\_uncertainty\_set}
  & $\ot$, $\wt$, conformal quantile $q_t$
  & Uncertainty set $\Ut$
  & Constructs $[\hat{x}_t - m_t,\;\hat{x}_t + m_t]$ with $m_t = q_t/(w_t + \epsilon)$.
\\[4pt]
\texttt{tighten\_action\_set}
  & $\Ut$, plant constraints
  & Safe action set $\Aset$
  & Shrinks the power/ramp feasible set by the margin $m_t$ on each SOC bound.
\\[4pt]
\texttt{repair\_action}
  & Candidate $\astar$, $\Aset$
  & Safe action $\asafe$
  & Projects $\astar$ onto $\Aset$ via constrained QP; returns $\astar$ if already feasible.
\\[4pt]
\texttt{emit\_certificate}
  & $\wt$, $m_t$, $\asafe$, prev.\ hash
  & Certificate record
  & Emits a hash-chained, per-step certificate with all DC3S fields for audit.
\\
\bottomrule
\end{tabular}
\end{table}

\section{Method Contracts}

Each method satisfies the following contracts:

\begin{description}[style=nextline]
  \item[\texttt{ingest\_telemetry}]
    \textbf{Pre:} raw packet is a byte-stream or structured dict.
    \textbf{Post:} $\ot \in \R^n$ with all NaN/missing values imputed
    by zero-order hold.
    \textbf{Latency:} $< 0.1$\,ms.

  \item[\texttt{compute\_oqe}]
    \textbf{Pre:} $\ot$ is well-formed; history buffer has $\ge 1$ entry.
    \textbf{Post:} $\wt \in [w_{\min}, 1.0]$; monotone in data quality.
    \textbf{Latency:} $< 0.5$\,ms.

  \item[\texttt{build\_uncertainty\_set}]
    \textbf{Pre:} $\wt > 0$, $q_t > 0$.
    \textbf{Post:} $\Ut$ is a non-empty interval; $m_t \le m_{\max}$.
    \textbf{Latency:} $< 0.1$\,ms.

  \item[\texttt{tighten\_action\_set}]
    \textbf{Pre:} $\Ut$ is non-empty.
    \textbf{Post:} $\Aset \subseteq \mathcal{U}$; $\Aset \neq \emptyset$
    when $\phi(x_t) = 1$ (Assumption~A3).
    \textbf{Latency:} $< 0.2$\,ms.

  \item[\texttt{repair\_action}]
    \textbf{Pre:} $\Aset \neq \emptyset$.
    \textbf{Post:} $\asafe \in \Aset$;
    $\|\asafe - \astar\|$ is minimised (projection).
    \textbf{Latency:} $< 1.0$\,ms.

  \item[\texttt{emit\_certificate}]
    \textbf{Pre:} all upstream fields computed.
    \textbf{Post:} certificate contains fields
    \{\texttt{step}, \texttt{w\_t}, \texttt{margin}, \texttt{action},
    \texttt{intervened}, \texttt{hash}, \texttt{prev\_hash}\}.
    \textbf{Latency:} $< 0.3$\,ms.
\end{description}

\section{Pipeline Diagram}

The six methods form a strict sequential pipeline within each dispatch
step:

\begin{center}
\begin{tikzpicture}[
  box/.style={draw, rounded corners, minimum width=2.2cm, minimum height=0.7cm,
              font=\small\ttfamily, fill=blue!8},
  arr/.style={-{Stealth[length=5pt]}, thick},
  node distance=0.35cm
]
\node[box] (a) {ingest};
\node[box, right=of a] (b) {oqe};
\node[box, right=of b] (c) {uncert};
\node[box, right=of c] (d) {tighten};
\node[box, right=of d] (e) {repair};
\node[box, right=of e] (f) {certify};
\draw[arr] (a) -- (b);
\draw[arr] (b) -- (c);
\draw[arr] (c) -- (d);
\draw[arr] (d) -- (e);
\draw[arr] (e) -- (f);
\end{tikzpicture}
\end{center}

\section{How to Implement a New Domain Adapter}

To extend ORIUS to a new CPS domain (e.g.\ HVAC, water treatment):

\begin{enumerate}
  \item \textbf{Subclass} \texttt{DomainAdapter} and implement the six
    methods above with domain-specific parsing, quality scoring, and
    constraint definitions.
  \item \textbf{Define the safety predicate} $\phi(x)$ for the new
    domain (e.g.\ temperature bounds for HVAC, chemical concentration
    for water).
  \item \textbf{Specify plant dynamics} $f(x, a)$ so that
    \texttt{tighten\_action\_set} can compute forward-reachable states.
  \item \textbf{Calibrate the conformal module} on domain-specific
    forecast residuals to obtain $q_t$.
  \item \textbf{Configure} \texttt{configs/dc3s.yaml} with
    domain-appropriate thresholds ($w_{\min}$, $m_{\max}$, $\alpha$).
  \item \textbf{Register} the adapter in the CPSBench runner so that
    the evaluation harness discovers it automatically.
\end{enumerate}

\noindent
The battery adapter in
\texttt{src/orius/dc3s/} serves as the canonical reference.
All six methods, their unit tests, and integration tests are documented
in the code-level API reference
(plan document \texttt{orius-plan/04-core-apis.md}).
